% kaobook-config.tex — кастомизация класса kaobook

%% ── Заголовки и нумерация ────────────────────────────────────────────────────
\setcounter{secnumdepth}{1}   % нумеруются главы и разделы, §§ — нет
\setcounter{tocdepth}{1}      % в оглавлении: главы и разделы

%% ── Боковые заметки ─────────────────────────────────────────────────────────
% Ширина основного текста + поля настраиваются в kaobook через geometry.
% Дополнительные настройки маргиналий:
\renewcommand{\marginfont}{\footnotesize\sffamily}

%% ── Колонтитулы ──────────────────────────────────────────────────────────────
\makeatletter
\renewcommand{\chaptermark}[1]{\markboth{\thechapter.\ #1}{}}
\renewcommand{\sectionmark}[1]{\markright{\thesection\ #1}}
\makeatother

%% ── Стиль листингов кода ─────────────────────────────────────────────────────
\definecolor{lst-bg}{HTML}{F6F8FA}
\definecolor{lst-comment}{HTML}{6A737D}
\definecolor{lst-keyword}{HTML}{D73A49}
\definecolor{lst-string}{HTML}{032F62}
\definecolor{lst-number}{HTML}{005CC5}

\lstset{
  backgroundcolor = \color{lst-bg},
  basicstyle      = \ttfamily\small,
  commentstyle    = \color{lst-comment}\itshape,
  keywordstyle    = \color{lst-keyword}\bfseries,
  stringstyle     = \color{lst-string},
  numberstyle     = \tiny\color{lst-number},
  breaklines      = true,
  frame           = single,
  framerule       = 0pt,
  rulecolor       = \color{lst-bg},
  numbers         = left,
  numbersep       = 8pt,
  showstringspaces= false,
  tabsize         = 2,
  captionpos      = b,
  xleftmargin     = 1em,
  xrightmargin    = 0em,
}

%% ── Стиль таблиц ─────────────────────────────────────────────────────────────
\renewcommand{\arraystretch}{1.25}

%% ── Цвета акцентов ───────────────────────────────────────────────────────────
% kaobook использует цвета из xcolor. Переопределяем под фирменный стиль:
\colorlet{kao-accent}{RoyalBlue!80!black}
